\documentclass[10pt,a4paper]{article}
\usepackage[utf8]{inputenc}
\usepackage[a4paper,left=3cm,right=2cm]{geometry}
\usepackage{amsmath}
\usepackage[thmmarks, amsmath, thref]{ntheorem}
\usepackage{amsfonts}
\usepackage{amssymb}
\usepackage{fancyhdr}
\pagestyle{fancy}
\usepackage{anysize}
\usepackage{graphicx}
\usepackage{float}
\usepackage{hyperref}
\usepackage{multirow}
\usepackage{enumitem}
\newtheorem{theorem}{Theorem}
\theoremsymbol{\ensuremath{\blacksquare}}
\newtheorem*{solution}{Solución:}

\usepackage{xcolor}

% Margenes y formalidades

\textwidth 6.25in % Margen Derecho y ancho del texto
\textheight 8.5in
\oddsidemargin 0in % Margen Izquierdo
\headheight 0in % Largo del texto
\leftmargin 1in
\parindent 0em % Salto entre la indentación (Sangría)https://www.overleaf.com/4519127669sdpptkhvsbfn
\topmargin -0.5in % Margen Superior
\parskip 0.5ex % Salto entre parrafos
\baselineskip 1pt

%\lhead{\includegraphics[scale=0.2]{logo.png}} % texto izquierda de la cabecera
%\chead{}                                      % texto centro de la cabecera
%\rhead{\includegraphics[scale=0.2]{FIG/logo-mba-2016.png}} % número de página a la derecha

\renewcommand{\headrulewidth}{0.4pt} % grosor de la línea de la cabecera
\renewcommand{\footrulewidth}{0.4pt} % grosor de la línea del pie

\begin{document}



\pagebreak
\vspace*{4mm} 

\begin{center}
\begin{huge}
\textbf{\\Ejercicios Econometría 3}
\end{Large}\\
%{\large 2023-2}\\
\end{center}


\section{Intervalos de Confianza: Proporción}

En una encuesta realizada a 500 personas, se encontró que 220 estaban de acuerdo con una nueva política pública. Calcula los intervalos de confianza para la proporción de personas que están de acuerdo con la política pública al 90\%, 95\% y 99\% de confianza.



\section{Intervalos de Confianza: Diferencia de medias}

Una empresa desea comparar los tiempos promedio de producción de dos plantas, Planta A y Planta B. Se toman muestras aleatorias de los tiempos de producción (en horas) de ambas plantas.

\begin{itemize}
    \item Para la \textbf{Planta A}, se selecciona una muestra de 50 observaciones con una media muestral de \( \bar{X}_A = 5,2 \) horas y una varianza conocida de \( \sigma_A^2 = 0,16 \) horas.
    \item Para la \textbf{Planta B}, se selecciona una muestra de 60 observaciones con una media muestral de \( \bar{X}_B = 4,9 \) horas y una varianza conocida de \( \sigma_B^2 = 0,25 \) horas.
\end{itemize}

 Verifique si la diferencia entre los tiempos de producción de las dos plantas es estadísticamente significativa.




\section{Regresión Lineal Multiple}

Se cuenta con la siguiente información para 5 individuos, donde $Y$ es la variable dependiente y $X_1$, $X_2$ son variables explicativas:

\[
\begin{array}{c|ccc}
\hline
\text{Individuo} & X_1 & X_2 & Y \\
\hline
1 & 4 & 1 & 5 \\
2 & 3 & 2 & 7 \\
3 & 5 & 3 & 10 \\
4 & 7 & 4 & 13 \\
5 & 6 & 4 & 12 \\
\hline
\end{array}
\]

\begin{itemize}
  \item \textbf{(Modelo de Regresión Lineal Simple)} Estime el modelo $Y = \beta_0 + \beta_1\,X_1 + \varepsilon$, ignorando la variable $X_2$. Calcule los coeficientes, el $R^2$, el $R^2$ ajustado y las sumas de cuadrados (SST, SSR y SSE).
  \item \textbf{Discuta la posibilidad de sesgo} al omitir $X_2$.
  \item \textbf{(Modelo de Regresión Lineal Múltiple)} Estime el modelo $Y = \beta_0 + \beta_1 X_1 + \beta_2 X_2 + \varepsilon$, calcule el nuevo $R^2$ (y las sumas de cuadrados) y compare con el modelo simple.
\end{itemize}




\end{document}